% Source: https://artofproblemsolving.com/community/c3h1571025_circle_geometry_problem
% Stable identity: aops-topic-1571025-q01
Place $A=(0,0)$, $B=(4,0)$, $C=(4,4)$, $D=(0,4)$. The semicircle has centre $O=(0,2)$ and radius $2$; the quarter-circle has centre $B$ and radius $4$.

\textbf{Intersection points.} The circles are $x^2+y^2-4y=0$ and $x^2+y^2-8x=0$. Subtracting gives $y=2x$, and then $5x^2-8x=0$, so the circles meet at $A=(0,0)$ and $X=\left(\dfrac{8}{5},\dfrac{16}{5}\right)$.

\begin{asy}
usepackage("amsmath");
size(5.5cm);
defaultpen(fontsize(10pt));
import graph;
pair A=(0,0),B=(4,0),C=(4,4),D=(0,4),O1=(0,2),X=(8/5,16/5);
real al=degrees(X-O1), be=degrees(X-B);
path semi=arc(O1,2,-90,90);
path quarter=arc(B,4,90,180);
path lens=arc(O1,2,-90,al)--arc(B,4,be,180)--cycle;
fill(lens,paleyellow);
draw(A--B--C--D--cycle,linewidth(0.9));
draw(semi,blue+linewidth(0.9));draw(quarter,red+linewidth(0.9));
draw(A--X,dashed);draw(O1--A,gray);draw(O1--X,gray);draw(B--X,gray);
dot(X);dot(O1);
label("$A$",A,SW);label("$B$",B,SE);label("$C$",C,NE);label("$D$",D,NW);label("$X$",X+0.12*dir(98),dir(98));label("$O$",O1,W);
label("$4$",(A+B)/2,S);
label("\mbox{}",(min(currentpicture,user=true)+max(currentpicture,user=true))/2,black); // keeps the text after the figure black when LaTeX typesets the labels inline
\end{asy}

The common region (shaded) is a lens cut by the chord $AX$ into two circular segments.

\textbf{Segment of the quarter-circle.} With $\beta=\angle ABX$: $\overrightarrow{BA}=(-4,0)$ and $\overrightarrow{BX}=\left(-\dfrac{12}{5},\dfrac{16}{5}\right)$, so $\cos\beta=\dfrac{3}{5}$ and $\sin\beta=\dfrac{4}{5}$. The segment has area
\[
\dfrac{4^2}{2}(\beta-\sin\beta)=8\beta-\dfrac{32}{5}.
\]

\textbf{Segment of the semicircle.} With $\alpha=\angle AOX$: $\overrightarrow{OA}=(0,-2)$ and $\overrightarrow{OX}=\left(\dfrac{8}{5},\dfrac{6}{5}\right)$, so $\cos\alpha=-\dfrac{3}{5}$, $\sin\alpha=\dfrac{4}{5}$ and $\alpha=\pi-\beta$. This segment (the part of the half-disc on the side of $AX$ towards $B$) has area
\[
\dfrac{2^2}{2}(\alpha-\sin\alpha)=2(\pi-\beta)-\dfrac{8}{5} .
\]

\textbf{Total.}
\[
\text{Area}=8\beta-\dfrac{32}{5}+2\pi-2\beta-\dfrac{8}{5}=2\pi+6\beta-8,
\]
where $\beta=\arccos\dfrac{3}{5}=2\arctan\dfrac{1}{2}$ (since $\tan\dfrac{\beta}{2}=\dfrac{\sin\beta}{1+\cos\beta}=\dfrac{1}{2}$). Hence
\[
\text{Area}=\boxed{2\pi-8+12\arctan\dfrac{1}{2}}\approx3.847 .
\]

\emph{Remark.} The post claims the answer can be found without trigonometry, but the angle $\beta$ (with $\cos\beta=\dfrac{3}{5}$) is not a rational multiple of $\pi$, because the only rational values of the cosine at rational multiples of $\pi$ are $0,\pm\dfrac{1}{2},\pm1$. So the area cannot be written using $\pi$ and radicals alone; an inverse trigonometric value is unavoidable.
